\documentclass[../cours_tle.tex]{subfiles}
\begin{document}
\setcounter{chapter}{20}
\chapter{Expériences aléatoires et loi binomiale}
\label{ch:21}
\begin{objectif}
Spécialité. Construire le modèle avant d'utiliser une formule. Prérequis : dénombrement et probabilités de première.
\end{objectif}
\section{Rappels : univers, événements et conditionnement}
Un univers fini $\Omega$ rassemble les issues ; un événement est une partie de $\Omega$. Une probabilité attribue des masses positives ou nulles de somme $1$ aux issues.
$P(A\cup B)=P(A)+P(B)-P(A\cap B)$ ; $P(\overline A)=1-P(A)$.
Si $P(A)>0$, $P(B\mid A)=P(A\cap B)/P(A)$. Cette probabilité décrit une information nouvelle, pas une causalité.
Pour une partition $(A_i)$ avec $P(A_i)>0$, les parties disjointes $B\cap A_i$ donnent
$P(B)=\sum_iP(A_i)P(B\mid A_i)$.
Si $P(B)>0$, $P(A_i\mid B)=P(A_i)P(B\mid A_i)/P(B)$ (formule de Bayes).
\section{Indépendance et arbres}
$A$ et $B$ sont indépendants si $P(A\cap B)=P(A)P(B)$. Si $P(A)>0$, cela équivaut à $P(B\mid A)=P(B)$.
L'incompatibilité signifie $A\cap B=\varnothing$ ; deux événements incompatibles de probabilités non nulles ne sont pas indépendants.

Une famille d'événements $A_1,\ldots,A_n$ est \emph{mutuellement indépendante} si, pour toute partie non vide $J$ de $\{1,\ldots,n\}$,
$P(\bigcap_{j\in J}A_j)=\prod_{j\in J}P(A_j)$.
L'indépendance deux à deux ne vérifie que les parties de taille 2 et ne suffit pas pour trois événements ou plus (exercice 21.5).
On peut remplacer certains événements par leurs complémentaires : pour une intersection $B$ des autres,
$P(B\cap\overline A)=P(B)-P(B\cap A)=P(B)(1-P(A))$ ; répéter cet argument.

Dans une succession d'expériences indépendantes, la probabilité d'une liste d'issues est le produit de leurs probabilités. Un arbre représente les choix successifs ; on multiplie le long d'un chemin et on additionne des chemins distincts conduisant à l'événement.
Sans indépendance, les probabilités des branches sont conditionnelles et peuvent varier suivant le chemin.
\section{Du schéma de Bernoulli à sa loi}
Une variable de Bernoulli de paramètre $p\in[0,1]$ vaut $1$ avec probabilité $p$ et $0$ avec probabilité $1-p$.
Un schéma de Bernoulli répète $n$ épreuves identiques et mutuellement indépendantes.
\begin{theoreme}
Le nombre $X$ de succès suit la loi binomiale $\mathcal B(n,p)$ :
$P(X=k)=\binom nk p^k(1-p)^{n-k}$ pour $k=0,\ldots,n$.
\end{theoreme}
\begin{demonstration}
Pour un choix fixé de $k$ positions de succès, l'indépendance donne $p^k(1-p)^{n-k}$. Il y a $\binom nk$ choix de positions, correspondant à des issues disjointes.
La somme des probabilités vaut $(p+(1-p))^n=1$ par le binôme.
Aux extrêmes $p=0$ ou $1$, utiliser directement les lois constantes pour éviter les ambiguïtés de notation $0^0$.
\end{demonstration}
Un tirage sans remise ne relève généralement pas de cette loi. Les probabilités de succès évoluent avec les tirages et les résultats sont dépendants.
\section{Calculer des événements et des seuils}
Pour $X\sim\mathcal B(5,0{,}2)$, $P(X=2)=10\times0{,}2^2\times0{,}8^3=0{,}2048$.
$P(X\ge1)=1-0{,}8^5=0{,}67232$.
Pour $0\le a\le b\le n$ entiers,
$P(a\le X\le b)=\sum_{k=a}^bP(X=k)$ ; avec la fonction cumulative $F(k)=P(X\le k)$, cela vaut $F(b)-F(a-1)$.

Le plus petit $k$ tel que $P(X>k)\le\alpha$ est le premier $k$ avec $F(k)\ge1-\alpha$. Il faut distinguer $>$ et $\ge$ : $P(X\ge k)=1-F(k-1)$.
Dans une surréservation, si chaque client se présente indépendamment avec probabilité $p$ et si $n$ réservations ont été prises pour $c$ places, le risque de dépassement est $P(X>c)$. Ce modèle dépend de l'hypothèse d'indépendance des présences.
\section{Concevoir une expérience et une simulation}
Une planche de Galton idéale ajoute $n$ choix droite/gauche indépendants de probabilité $1/2$. Le numéro du compartiment suit $\mathcal B(n,1/2)$.
Simuler un échantillon consiste à répéter ce mécanisme ; les fréquences obtenues fluctuent autour des probabilités théoriques. Une simulation n'établit ni l'indépendance réelle du dispositif ni une égalité exacte de probabilités.
Pour comparer plusieurs valeurs de $n$, calculer le risque pour chacune et justifier le choix minimal ou maximal sous la contrainte, sans seulement afficher une valeur trouvée.
\begin{finchapitre}
\section{Utilité et suite du parcours}
Les lois discrètes sont la base de l'inférence statistique. La factorisation des probabilités prépare les modèles produits et les chaînes de Markov, où l'indépendance sera remplacée par une dépendance contrôlée.
\begin{retenir}
Vérifier nombre d'épreuves fixé, même probabilité et indépendance. Traduire précisément « au moins », « au plus » et « strictement plus ».
\end{retenir}
\end{finchapitre}
\end{document}
